\documentclass[12pt,a4paper]{amsart}

\usepackage{amsmath,amssymb,amsthm}
\usepackage{mathtools}
\usepackage[ngerman]{babel}
\usepackage{hyperref}
\usepackage{geometry}
\geometry{margin=2.5cm}

% -------------------------------------------------------
% Theorem-Umgebungen
% -------------------------------------------------------
\newtheorem{theorem}{Satz}[section]
\newtheorem{lemma}[theorem]{Lemma}
\newtheorem{corollary}[theorem]{Korollar}
\newtheorem{proposition}[theorem]{Proposition}
\theoremstyle{definition}
\newtheorem{definition}[theorem]{Definition}
\newtheorem{remark}[theorem]{Bemerkung}
\newtheorem{conjecture}[theorem]{Vermutung}

% -------------------------------------------------------
% Eigene Befehle
% -------------------------------------------------------
\newcommand{\N}{\mathbb{N}}
\newcommand{\Z}{\mathbb{Z}}
\newcommand{\R}{\mathbb{R}}
\newcommand{\C}{\mathbb{C}}
\DeclareMathOperator{\Re}{Re}
\DeclareMathOperator{\Im}{Im}
\DeclareMathOperator{\tr}{Spur}
\DeclareMathOperator{\spec}{Spec}

% -------------------------------------------------------
\begin{document}
% -------------------------------------------------------

\title{Ans\"atze zur Riemann-Hypothese:\\
       Hilbert--P\'olya, Berry--Keating, Zufallsmatrizen\\
       und bekannte Barrieren}

\author{Michael Fuhrmann}
\address{Specialist Mathematics Project, Build 84}
\email{specialist-maths@localhost}
\date{\today}

\begin{abstract}
Die Riemann-Hypothese (RH) -- alle nicht-trivialen Nullstellen von $\zeta(s)$
liegen auf der kritischen Geraden $\Re(s) = \tfrac{1}{2}$ -- widersteht
seit 1859 jedem Beweisversuch.
Wir beschreiben die wichtigsten Ans\"atze:
die \emph{Hilbert--P\'olya-Vermutung} (Nullstellen = Eigenwerte eines
selbstadjungierten Operators),
Berry und Keatings \emph{semiklassische Quantisierung} des klassischen
Hamiltonians $H = xp$,
die von Montgomery (1973) entdeckte und von Odlyzko numerisch verifi\-zierte
\emph{GUE-Verbindung} und den tiefen Zusammenhang mit
\emph{Quantenchaos}.
Au\ss erdem werden die bekannten Barrieren beschrieben sowie das
Teilresultat von Conrey (1989), dass mehr als $\tfrac{2}{5}$ aller
Nullstellen auf der kritischen Geraden liegen.
\end{abstract}

\maketitle
\tableofcontents

% ================================================================
\section{Einleitung: Aussage und Bedeutung}
% ================================================================

\begin{conjecture}[Riemann-Hypothese, 1859 \cite{Riemann1859}]
\label{conj:rh}
Alle nicht-trivialen Nullstellen $\rho$ der Riemann'schen Zeta-Funktion
erf\"ullen $\Re(\rho) = \tfrac{1}{2}$.
\end{conjecture}

Die Riemann-Hypothese ist das ber\"uhmteste ungel\"oste Problem der Mathematik
und eines der sieben Millennium-Preisaufgaben des Clay Mathematics Institute
(Preisgeld: 1.000.000~USD).

\begin{remark}[Was bekannt ist]
\begin{itemize}
  \item $\zeta(s) \ne 0$ f\"ur $\Re(s) \ge 1$ (Primzahlsatz, Hadamard 1896).
  \item $\zeta(s) \ne 0$ im nullstellenfreien Gebiet $\sigma > 1 - c/\log t$
    (de la Vall\'ee~Poussin 1899).
  \item Mehr als $\tfrac{2}{5}$ aller Nullstellen liegen auf der kritischen
    Geraden (Conrey 1989 \cite{Conrey1989}).
  \item \"Uber $10^{13}$ Nullstellen numerisch auf der Geraden verifiziert
    (Platt--Trudgian 2021).
  \item Alle bisher gefundenen Nullstellen erf\"ullen $\Re(\rho) = \tfrac{1}{2}$.
\end{itemize}
\end{remark}

% ================================================================
\section{Die Hilbert--P\'olya-Vermutung}
% ================================================================

\begin{conjecture}[Hilbert--P\'olya]
\label{conj:hilbert_polya}
Es existiert ein selbstadjungierter (hermitescher) Operator $H$ auf
einem Hilbertraum $\mathcal{H}$, dessen Eigenwerte genau die
Imagin\"arteile $\gamma_n$ der nicht-trivialen Nullstellen
$\rho_n = \tfrac{1}{2} + i\gamma_n$ sind:
\[
  \spec(H) = \{\gamma_n : \zeta(\tfrac{1}{2} + i\gamma_n) = 0,\; \gamma_n > 0\}.
\]
\end{conjecture}

\begin{remark}[Entstehung]
David Hilbert und George P\'olya schlugen unabh\"angig voneinander um
1910--1920 vor, dass RH aus der Existenz eines solchen hermiteschen Operators
folgen k\"onnte, da Eigenwerte selbstadjungierter Operatoren stets reell sind.
Falls ein solches $H$ existiert, sind alle $\gamma_n \in \R$, was
\"aquivalent zu RH ist.
\end{remark}

\begin{proposition}[Reduktion auf RH]
Wenn ein selbstadjungierter Operator $H$ mit Eigenwerten $\{\gamma_n\}$
(den Imagin\"arteilen der nicht-trivialen Nullstellen) existiert,
dann ist RH wahr.
\end{proposition}

\begin{proof}
Selbstadjungierte Operatoren haben reelles Spektrum. Ist $\gamma_n \in \R$
f\"ur alle $n$, so gilt $\rho_n = \tfrac{1}{2} + i\gamma_n$ mit
$\gamma_n \in \R$, also $\Re(\rho_n) = \tfrac{1}{2}$.
\end{proof}

\begin{remark}[Status]
Kein explizites $H$ gem\"a\ss\ Vermutung~\ref{conj:hilbert_polya} wurde bisher
gefunden. Die Schwierigkeit liegt in der nat\"urlichen Konstruktion von $H$
mit dem richtigen Spektrum.
Berry--Keating (Abschnitt~\ref{sec:berry_keating}) ist der am weitesten
entwickelte Ansatz.
\end{remark}

% ================================================================
\section{Berry--Keating: Der \texorpdfstring{$xp$}{xp}-Hamiltonian}
\label{sec:berry_keating}
% ================================================================

\begin{definition}[Klassischer $xp$-Hamiltonian]
\label{def:xp}
Der \emph{Berry--Keating-Hamiltonian} ist der klassische Hamiltonian
\[
  H = xp,
\]
wobei $x$ der Ort und $p$ der Impuls ist, wirkend auf den klassischen
Phasenraum $\R_{>0} \times \R_{>0}$.
\end{definition}

\begin{theorem}[Semiklassisches Spektrum von $H=xp$, Berry--Keating 1999]
\label{thm:berry_keating}
Berry und Keating \cite{BerryKeating1999} beobachteten:
\begin{enumerate}
  \item[\textup{(i)}] Die klassischen Trajektorien von $H = xp$ sind
    die Hyperbeln $xp = E$ (konstantes $E > 0$) im Phasenraum.
  \item[\textup{(ii)}] Die WKB-Quantisierungsbedingung (Bohr--Sommerfeld)
    f\"ur Orbits $\{(x,p) : xp = E,\; a \le x \le b\}$ liefert:
    \[
      E_n \approx \frac{2\pi n}{\log(n/2\pi e)} \quad (n \to \infty),
    \]
    was mit der Asymptotik $\gamma_n \approx 2\pi n/\log(n/2\pi)$ der
    Nullstellen \"ubereinstimmt.
  \item[\textup{(iii)}] Der quantisierte Hamiltonian
    $\hat H = \tfrac{1}{2}(x\hat p + \hat p x)$
    (symmetrisiert) auf $L^2(\R_{>0})$ ist formal selbstadjungiert,
    aber hat auf einem nat\"urlichen Gebiet nicht genau die $\{\gamma_n\}$
    als Eigenwerte.
\end{enumerate}
\end{theorem}

\begin{proof}[Diskussion]
Der klassische Phasenraum $xp = E$ hat Fl\"ache (in $\log x, \log p$-Koordinaten)
$\log E$.
Die Anzahl der Zust\"ande bis zur Energie $E$ ist somit
$N(E) \approx E/(2\pi) \cdot \log(E/2\pi e) + 7/8$,
was mit der Riemann--von-Mangoldt-Formel aus Paper~22 \"ubereinstimmt.

Das fehlende Element ist eine Randbedingung oder Regularisierung,
die den formalen Operator in einen echten selbstadjungierten Operator mit
genau den $\gamma_n$ als Eigenwerten verwandelt.
Verschiedene Vorschl\"age wurden untersucht, keiner liefert bisher
einen vollst\"andigen rigoros Beweis.
\end{proof}

\begin{remark}[Connes' Umformulierung]
Alain Connes \cite{Connes1999} schlug eine Umformulierung mit
\emph{nichtkommutativer Geometrie} vor: Die Zeta-Nullstellen sollten
die \glqq fehlenden Eigenwerte\grqq{} eines nat\"urlichen
Dirac-artigen Operators auf einem nichtkommutativen Raum sein.
Dieses Programm ist tief, aber noch unbewiesen.
\end{remark}

% ================================================================
\section{Zufallsmatrizen und GUE-Statistik}
\label{sec:gue}
% ================================================================

\begin{definition}[Gaussian Unitary Ensemble]
\label{def:gue}
Das \emph{Gaussian Unitary Ensemble} (GUE) der Zufallsmatrizentheorie
ist die Wahrscheinlichkeitsverteilung auf $n \times n$ hermiteschen Matrizen $M$
mit Dichte proportional zu $e^{-\tr M^2}$.

Die \emph{GUE-Paarkorrelationsfunktion} (f\"ur normalisierte Abst\"ande) ist:
\[
  r_2(s) \;=\; 1 - \left(\frac{\sin(\pi s)}{\pi s}\right)^2.
\]
\end{definition}

\begin{theorem}[Montgomerys Paarkorrelation \cite{Montgomery1973}]
\label{thm:montgomery_gue}
Unter Annahme der Riemann-Hypothese gilt f\"ur glatte Testfunktionen $f$:
\[
  \frac{1}{N(T)} \sum_{\gamma, \gamma'} f\!\left(\frac{(\gamma-\gamma')\log T}{2\pi}\right)
  \;\xrightarrow{T \to \infty}\;
  \int_{-\infty}^{\infty} f(x)\,r_2(x)\,dx.
\]
Dies ist genau die GUE-Paarkorrelationsfunktion $r_2(s) = 1-(\sin\pi s/\pi s)^2$.
\end{theorem}

\begin{remark}
Montgomery bewies dies f\"ur Testfunktionen mit $\hat f$ tr\"ager in $(-1,1)$.
Die volle Vermutung (alle $f$) ist unbewiesen.
Die Aussage wurde von Freeman Dyson (bei einem Tee am IAS 1972) als
GUE-Verteilung aus der Kernphysik identifiziert.
\end{remark}

\begin{theorem}[Odlyzkos numerische Verifikation]
\label{thm:odlyzko_gue}
Odlyzko \cite{Odlyzko1987} berechnete Statistiken der Nullstellen in der
N\"ahe der $10^{20}$-ten Nullstelle und verglich mit GUE-Vorhersagen:
\begin{itemize}
  \item N\"achste-Nachbar-Abstandsverteilung: Abweichung $< 1\,\%$ von GUE.
  \item $k$-Niveau-Korrelationsfunktionen: \"Ubereinstimmung f\"ur $k=2,3,4$.
  \item Niveau-Varianz $\Sigma^2(L)$: \"Ubereinstimmung mit GUE.
\end{itemize}
Dies gilt als \"uberzeugendstes numerisches Indiz f\"ur RH und
die Hilbert--P\'olya-Vermutung.
\end{theorem}

% ================================================================
\section{Quantenchaos und die Selberg-Spurformel}
% ================================================================

\begin{remark}[Verbindung zum Quantenchaos]
Die GUE-Universalit\"atsklasse tritt bei Quantensystemen auf, deren
klassisches Pendant \emph{chaotisch} ist (Berry--Tabor-Vermutung 1977,
Bohigas--Giannoni--Schmit-Vermutung 1984). Dies legt nahe:

\textbf{Der hypothetische Operator $H$ (Hilbert--P\'olya) sollte ein
klassisch chaotisches System quantisieren.}

Indizien:
\begin{itemize}
  \item Nullstellen von $\zeta$ haben GUE-Statistik.
  \item Eigenwerte quantisierter chaotischer Systeme (z.\,B.\ Geod\"aten auf
    hyperbolischen Fl\"achen) haben ebenfalls GUE-Statistik.
  \item Die Selberg-Spurformel f\"ur hyperbolische Fl\"achen ist analog
    zur expliziten Formel f\"ur $\zeta$.
\end{itemize}
\end{remark}

\begin{theorem}[Selberg-Spurformel (f\"ur kompakte hyperbolische Fl\"achen)]
\label{thm:selberg_trace}
Sei $\mathcal{M}$ eine kompakte hyperbolische Fl\"ache mit Geod\"atenl\"angen
$l(\gamma)$. Das Spektrum des Laplace-Operators
$\{0 = \lambda_0 < \lambda_1 \le \lambda_2 \le \ldots\}$,
$\lambda_n = \tfrac{1}{4} + r_n^2$, erf\"ullt:
\[
  \sum_n h(r_n) \;=\; \frac{\text{Fl}(\mathcal{M})}{4\pi}
  \int_\R h(r)\,r\tanh(\pi r)\,dr
  \;+\; \sum_{\substack{\text{prim.}\\ \text{Geod.}\ \gamma}} \sum_{k=1}^\infty
  \frac{l(\gamma)}{2\sinh(k l(\gamma)/2)}\,\hat h(k l(\gamma)).
\]
\end{theorem}

\begin{remark}
Die Selberg-Spurformel ist das \emph{geometrische Analogon} der
expliziten Formel:
Eigenwerte (Spektralseite) werden mit Geod\"atenl\"angen (geometrische Seite)
verbunden, genau wie Nullstellen von $\zeta$ mit Primzahlen.
Diese Analogie legt dringend nahe, dass die Nullstellen von $\zeta$
das Spektrum eines Operators sind -- aber welches Operators, ist offen.
\end{remark}

% ================================================================
\section{Teilresultate und bekannte Absch\"atzungen}
% ================================================================

\begin{theorem}[Nullstellenfreies Gebiet]
\label{thm:zero_free}
Es gibt eine positive Konstante $c > 0$, sodass $\zeta(s) \ne 0$ gilt f\"ur
\[
  \sigma > 1 - \frac{c}{\log(|t|+2)}.
\]
Unter der Riemann-Hypothese erstreckt sich dies auf $\sigma > \tfrac{1}{2}$.
\end{theorem}

\begin{theorem}[Levinson--Conrey: Anteil auf der kritischen Geraden]
\label{thm:levinson_conrey}
\begin{itemize}
  \item Levinson (1974): Mindestens $\tfrac{1}{3}$ aller nicht-trivialen
    Nullstellen liegen auf der kritischen Geraden.
  \item Conrey (1989) \cite{Conrey1989}: Mindestens $\tfrac{2}{5}$ aller
    nicht-trivialen Nullstellen liegen auf der kritischen Geraden.
  \item Feng (2011): Mindestens $0{,}41$ der Nullstellen liegen auf der Geraden.
\end{itemize}
\end{theorem}

\begin{proof}[Beweisskizze f\"ur Levinsons $\tfrac{1}{3}$]
Levinsons Methode z\"ahlt Nullstellen auf der kritischen Geraden
durch Untersuchung der \glqq mollifiierten\grqq{} Funktion
$\zeta(\tfrac{1}{2}+it) \cdot M(\tfrac{1}{2}+it)$,
wobei $M$ ein Dirichlet-Polynom ist.
Vorzeichenwechsel des Realteils dieses Produkts zeigen Nullstellen auf
der Geraden an.
Der Anteil $\tfrac{1}{3}$ folgt aus sorgf\"altiger Absch\"atzung der
Momente des mollifiierten $\zeta$.
\end{proof}

% ================================================================
\section{Barrieren und Hindernisse}
% ================================================================

\begin{remark}[Warum RH so schwer ist: bekannte Barrieren]
Mehrere Barrieren verhindern, dass naive Ans\"atze zum Erfolg f\"uhren:

\textbf{Barriere 1: Explizite Formeln reichen nicht.}
Die explizite Formel dr\"uckt $\pi(x)$ durch Nullstellen aus, aber um
daraus auf die Lage der Nullstellen zu schlie\ss en, braucht man
unabh\"angige Information.

\textbf{Barriere 2: Nullstellenfreie-Gebiet-Methoden.}
Klassische Methoden geben nullstellenfreie Gebiete $\sigma > 1 - c/\log t$,
k\"onnen aber ohne neue Ideen die Grenze $\sigma > \tfrac{1}{2}$ nicht
erreichen.
Siegel-Nullstellen (Nullstellen sehr nahe an $\sigma = 1$ f\"ur
$L$-Funktionen) sind ein gro\ss es Hindernis auch im GRH-Kontext.

\textbf{Barriere 3: Kein bekannter Operator.}
Der Hilbert--P\'olya-Ansatz erfordert einen expliziten selbstadjungierten
Operator. Alle bisherigen Kandidaten sind nicht vollst\"andig rigoros.

\textbf{Barriere 4: GUE-Evidenz ist kein Beweis.}
Montgomerys Satz ist bedingt (setzt RH voraus) und Odlyzkos Berechnungen
sind numerisch -- sie liefern keinen Beweisansatz.

\textbf{Barriere 5: Die Analogie mit Funktionenk\"orpern ist unvollst\"andig.}
Weil (1940) und Deligne (1974) bewiesen die Riemann-Hypothese f\"ur
Zeta-Funktionen \"uber endlichen K\"orpern (Weil-Vermutungen).
Die Methoden (algebraische Geometrie \"uber endlichen K\"orpern)
\"ubertragen sich jedoch nicht direkt.
\end{remark}

% ================================================================
\section{Verbindung zur Physik: Quantenchaos und dar\"uber hinaus}
% ================================================================

\begin{remark}[Dreifache Analogie]
Die folgende Analogie durchzieht alle RH-Ans\"atze:
\[
\begin{array}{ccc}
  \text{Zahlentheorie} & \longleftrightarrow & \text{Quantenmechanik} \\
  \hline
  \text{Primzahlen } p & \longleftrightarrow & \text{Periodische Orbits} \\
  \zeta(s) & \longleftrightarrow & \text{Spektraldeterminante} \\
  \text{Nullstellen } \gamma_n & \longleftrightarrow & \text{Energie-Eigenwerte} \\
  \text{Euler-Produkt} & \longleftrightarrow & \text{Spurformel} \\
  \text{GUE-Statistik} & \longleftrightarrow & \text{Quantenchaos}
\end{array}
\]
Diese Analogie ist sehr suggestiv, hat aber noch zu keinem Beweis gef\"uhrt.
\end{remark}

% ================================================================
\section{Schlussfolgerung}
% ================================================================

Die Riemann-Hypothese hat anderthalb Jahrhunderte tiefer Mathematik
inspiriert: die analytische Theorie von $\zeta$ und $L$-Funktionen,
die Verbindung zur Quantenmechanik (Hilbert--P\'olya, Berry--Keating),
die tiefe GUE-Statistik (Montgomery, Odlyzko) und die Selberg-Spurformel.

Trotzdem bleibt ein Beweis unerreichbar. Die Barrieren sind bekannt:
klassische Methoden k\"onnen $\sigma = \tfrac{1}{2}$ nicht unterschreiten;
kein selbstadjungierter Operator wurde explizit gefunden;
die Deligne-Methode aus der algebraischen Geometrie \"ubertr\"agt sich nicht.
Ben\"otigt wird eine genuinely neue Idee.

% ================================================================
\begin{thebibliography}{10}

\bibitem{Riemann1859}
B.~Riemann,
\emph{\"Uber die Anzahl der Primzahlen unter einer gegebenen Gr\"o\ss e},
Monatsberichte der Berliner Akademie (1859), 671--680.

\bibitem{Montgomery1973}
H.~L.~Montgomery,
The pair correlation of zeros of the zeta function,
in: \emph{Analytic Number Theory}, Proc.\ Sympos.\ Pure Math.~XXIV,
Amer.\ Math.\ Soc., 1973, S.~181--193.

\bibitem{Odlyzko1987}
A.~M.~Odlyzko,
On the distribution of spacings between zeros of the zeta function,
\emph{Math.\ Comp.}\ \textbf{48} (1987), 273--308.

\bibitem{BerryKeating1999}
M.~V.~Berry und J.~P.~Keating,
The Riemann zeros and eigenvalue asymptotics,
\emph{SIAM Rev.}\ \textbf{41} (1999), 236--266.

\bibitem{Connes1999}
A.~Connes,
Trace formula in noncommutative geometry and the zeros of the Riemann
zeta function,
\emph{Selecta Math.\ (N.S.)}\ \textbf{5} (1999), 29--106.

\bibitem{Conrey1989}
J.~B.~Conrey,
More than two fifths of the zeros of the Riemann zeta function are on
the critical line,
\emph{J.\ Reine Angew.\ Math.}\ \textbf{399} (1989), 1--26.

\bibitem{Davenport2000}
H.~Davenport,
\emph{Multiplicative Number Theory}, 3.~Aufl.,
Springer, New York, 2000.

\bibitem{Titchmarsh1986}
E.~C.~Titchmarsh,
\emph{The Theory of the Riemann Zeta-Function}, 2.~Aufl.,
Oxford University Press, 1986.

\end{thebibliography}

\end{document}
